\begin{resumo}
  A qualidade de software constitui elemento central na Engenharia de Software, sendo a manutenibilidade uma de suas características mais relevantes para a evolução e a sustentabilidade de sistemas ao longo do tempo. Ferramentas tradicionais de análise estática, contudo, apresentam limitações na captura de aspectos semânticos e contextuais do código, lacuna que motiva a investigação do uso de Modelos de Linguagem de Grande Escala (LLMs) como alternativa complementar de avaliação. Nesse cenário, a forma como as instruções são formuladas — os padrões de engenharia de prompt — surge como fator potencialmente relevante para os resultados produzidos por esses modelos. Este trabalho tem como objetivo investigar como diferentes padrões de engenharia de prompt influenciam a aferição do indicador de manutenibilidade de produtos de software por LLMs, tomando como referência modelos de qualidade consolidados na literatura. A pesquisa foi estruturada a partir do método Goal-Question-Metric (GQM) para definição da questão de pesquisa, seguida de uma revisão estruturada da literatura, conduzida com apoio do protocolo PICO na base Scopus, da qual resultaram 38 estudos analisados. A síntese dos achados evidenciou que a manutenibilidade é predominantemente operacionalizada por meio de medidas como complexidade, code smells, duplicação e dívida técnica. Além disso, constatou-se que não há um padrão de prompt específico consolidado para a manutenibilidade, embora estratégias que reduzem a janela de contexto de análise e dividem as solicitações em tarefas menores tendam a favorecer o desempenho das LLMs. A partir dessas evidências, foi planejado um experimento controlado, estruturado segundo um Randomized Complete Block Design (RCBD), no qual três padrões de prompt (zero-shot, few-shot e chain-of-thought) serão aplicados a arquivos de código-fonte de projetos open source do domínio de aplicações web, utilizando três ferramentas de LLM (ChatGPT, Claude e DeepSeek), tendo como variáveis de bloqueio dois modelos de qualidade de referência (Q-Rapids e MeasureSoftGram) e utilizando o método SQALE (via SonarQube) para triangulação complementar dos resultados. Espera-se, com a execução desse experimento na etapa subsequente do trabalho, verificar se há diferenças significativas entre os padrões de prompt na aferição da manutenibilidade e em que medida os resultados obtidos por LLMs se alinham aos modelos de qualidade tradicionalmente utilizados na Engenharia de Software.

  \vspace{\onelineskip}
     
  \noindent
  \textbf{Palavras-chave}: Modelos de Linguagem de Grande Escala. Manutenibilidade de Software. Engenharia de Prompt. Experimento Controlado. Qualidade de Software.
\end{resumo}
